\subsection{Assurance-vie}
\label{sec:assurance_vie}

L'assurance-vie bénéficie d'un régime fiscal dérogatoire qui s'applique aussi bien pendant la phase d'épargne qu'au moment de la transmission \cite{cgi_990i, cgi_757b, bofip_mutations_av}.

\subsubsection{Fonctionnement et frais}
\label{sec:rendement_prelevements}

Le capital placé sur une assurance-vie évolue en fonction de la performance des supports choisis (fonds en euros, unités de compte). Cependant, cette performance est systématiquement diminuée par des \textbf{frais de gestion annuels} prélevés par l'assureur. Ces frais, généralement compris entre 0,5\% et 1\% de l'encours, réduisent d'autant la croissance du capital sur le long terme.

\subsubsection{Fiscalité des retraits (rachats)}
\label{sec:rachats_partiels}

Lors d'un retrait (appelé "rachat"), l'épargnant n'est pas imposé sur la totalité de la somme retirée, mais uniquement sur la \textbf{part de gains} que ce retrait contient. Cette part est calculée proportionnellement au poids des gains dans la valeur totale du contrat.

Après 8 ans, les gains bénéficient d'un \textbf{abattement annuel de 4 600 €} (9 200 € pour un couple). Seuls les gains dépassant ce seuil sont soumis à l'impôt sur le revenu. En revanche, les prélèvements sociaux (17,2\%) restent dus sur la totalité des gains retirés.

\paragraph{Exemple.} Pour un contrat de valeur \(120000 €\) avec \(100000 €\) de primes nettes versées, soit \(20000 €\) de gains latents, un retrait partiel de 10000 € induit un gain imposable de :
La part de gains dans ce retrait est de : \(10\,000 \times \frac{20\,000}{120\,000} = 1\,667 €\).
Si le contrat a plus de 8 ans, ces \(1 667 €\) sont entièrement couverts par l'abattement annuel. L'épargnant ne paie donc \textbf{aucun impôt sur le revenu}, mais il doit tout de même s'acquitter des prélèvements sociaux : \(1\,667 \times 17,2\% = 287 €\).

\subsubsection{Fiscalité en cas de transmission}
\label{sec:fiscalite_transmission}

Au décès du souscripteur, la transmission du capital est soumise à une fiscalité en deux temps :
\begin{enumerate}
    \item \textbf{Prélèvements sociaux (PS) :} L'assureur prélève d'abord 17,2\% sur la totalité des gains accumulés et non encore taxés. Le capital transmis aux bénéficiaires est donc net de ces PS.
    \item \textbf{Droits de succession spécifiques :} Le capital restant bénéficie ensuite d'un abattement de 152 500 € par bénéficiaire (pour les versements effectués avant 70 ans). La part excédant cet abattement est taxée à des taux forfaitaires \cite{bofip_mutations_av}:
    \begin{itemize}
        \item 20\% jusqu'à 700 000 € (après abattement),
        \item 31,25\% au-delà.
    \end{itemize}
\end{enumerate}
L'héritage net correspond donc au capital final, diminué des prélèvements sociaux et des droits de succession spécifiques à l'assurance-vie.
